\scanpage{69}
\mspage{D}
% archive foliation: 68
D'ailleurs les modèles propres et normaux, pour la relation de domination,
forment un ordonné \uline{filtrant croissant},
\deleted{\unsure{on a}}
Soit $U_0$ un \deleted{tel} modèle
\deleted{\unsure{de l'un des ens pts réguliers}}
\inserted{régulier de $K$},
et soit $X_{\unsure{1|\alpha}}$
\deleted{dominant $X_0$}
\inserted{propre \deleted{\unsure{normal}} \uline{dominant}}.

\margin{\diagram{\begin{tikzcd}
&X''&\\
\deleted{\unsure{U}}X\arrow[ur,no head]\arrow[r,no head]&X'\deleted{\unsure{U'}}&
\end{tikzcd}}}

\deleted{Soit $\unsure{U_1|U_0}$ \unsure{l'ens des pts} de $U_0$ où
l'application rationnelle évidente \deleted{\unsure{$X$}}
$U_0\to X_{\unsure{1|\alpha}}$ induit
\deleted{\unsure{par $f$}} est définie,
c'est \inserted{le compl. du lieu d'}
\unsure{un ens fermé} de codim $\geq2$,
\deleted{\unsure{donc d'après le}}
et on a donc
\deleted{complémentaire dans $X_0$ d'un ens fermé de codim $\geq2$.}
\unsure{On a un}
\deleted{\unsure{hom canonique}},
diagramme commutatif, donc on voit que on a un hom surjectif}

\margin{\diagram{\begin{tikzcd}
X_{\unsure{1|\alpha}}\arrow[d,"f"']&&\\
X_0\arrow[r,phantom,"\supset" description]&
U_0\arrow[r,phantom,"\supset" description]&
U_{\unsure{1|\alpha}}\arrow[ull,"g"']
\end{tikzcd}}}
\margin{\diagram{\begin{tikzcd}
\pi_1(X_{\unsure{1|\alpha}})\arrow[dd,"f_*"']&
\deleted{\pi_1(U_{\unsure{1|\alpha}})}\arrow[l]&\\
&&\\
\pi_1(X_0)&\pi_1(U_0)\arrow[l]&
\pi_1(U_1)\arrow[l,"\simeq" above,"\text{th. de pureté}" below]\arrow[uull,"g_*"]
\end{tikzcd}}}

L'application rationnelle évidente $U\to X$ est définie sur le
complémentaire d'une partie fermée de codim $\geq2$,
on voit que (par le th. de pureté) on en tire un hom
\uline{surjectif} (car $X$ normal)
\[
\pi_1(U)\longrightarrow\pi_1(X)
\]
d'où à la limite un hom \uuline{surjectif}
\[
\pi_1(U;\Omega)\longrightarrow
\varprojlim_{\substack{X\ \text{modèle}\\\text{propre}}}
\pi_1(X;\Omega)
\]
où $\Omega$ est une extension alg. close de $K$.

\margin{\deleted{\unsure{of}}}
